% Suggested rewrite of the residual-histogram paragraph in hi_disk_size_environments.tex.
% Inline numbers use the macros from autogen/macros_hcg_comparison.tex
% (\HCGSizeFractionPercent, \HCGSizeReductionPercent,
%  \CliffProbAmigaGreaterHCGPercent), so this paragraph stays in sync with
% the script outputs automatically. See products/hcg_residual_statistics_*.json
% for the underlying numbers.

Figure~\ref{fig:residuals_hist} shows the density-normalised distributions
of the residuals. The AMIGA histogram is broadly distributed and
centred at zero, as expected from the reference baseline, with comparable
weight on both the positive and negative sides.  The HCG histogram has a
similar overall shape but is systematically shifted toward negative
residuals, with only a small tail extending to positive values.  The two distributions have
comparable widths ($\sigma_{\rm obs}^{\rm HCG} \simeq 0.18$~dex versus
$\sigma_{\rm obs}^{\rm AMIGA} \simeq 0.16$~dex, a ratio of $1.13$), so the
difference is dominated by a shift of the mean rather than a difference
in width. The mean offset implies that, on average, HCG galaxies
have \HI\ discs about \HCGSizeFractionPercent\% of the size expected for
isolated galaxies of the same optical diameter (i.e.\ a
$\sim$\HCGSizeReductionPercent\% reduction).

To test whether the AMIGA and HCG residuals could plausibly be drawn from
the same parent distribution, we apply complementary two-sample tests
that make minimal assumptions about the data (Mann--Whitney $U$,
Kolmogorov--Smirnov, and Anderson--Darling). All three reject the null
at very high significance, as summarised in Table~\ref{table:stat_tests}.
Our Cliff's $\delta$ effect-size measure indicates that a randomly chosen AMIGA galaxy has a larger residual than a
randomly chosen HCG galaxy roughly \CliffProbAmigaGreaterHCGPercent\% of
the time.
